


The \emph{Heaviside step function} is an example of a function
with a jump discontinuity.
It is defined as follows:
\[
	H(x) \eqdef \begin{cases}
				\; 1, 	& \text{if } x \geq 0 \\
				\; 0, 	& \text{if } x < 0.
			\end{cases}
\]
The function is zero for negative values of $x$,
then suddenly jumps to one at $x=0$,
as shown in Figure~\ref{fig:limits_examples}~(c).
The limit as $x$ approaches $x=0$ from the left is $\lim_{x\to 0^-} H(x) = 0$.
The limit at $x=0$ from the right is $\lim_{x\to 0^+} H(x) = 1$.
The two limits are different,
$\lim_{x\to 0^-} H(x) = 0 \neq 1 = \lim_{x\to 0^+} H(x)$,
so the function is \emph{discontinuous} at $x=0$.



% EXAMPLE  Heaviside
\noindent
To calculate the limit from the left or the right of a number,
we provide a fourth argument \tt{"-"} or \tt{"+"} to \tt{sp.limit}.
%		The following SymPy calculations confirm the limits of the Heaviside function
%		when approaching $x=0$ from the left and from the right.


\begin{codeblock}[]
>>> from sympy import Heaviside
>>> sp.limit(Heaviside(x,1), x, 0, "-")
0
>>> sp.limit(Heaviside(x,1), x, 0, "+")
1
\end{codeblock}

%	\noindent
%	See Figure~\ref{fig:limits_examples}~(c) for an illustration.
